\section{Formalization Notes}\label{sec:formalization}

\subsection{Overview}

The function field module is part of a larger Lean~4 formalization of the
reduction of Mullin's Conjecture to open hypotheses, totaling approximately
72{,}000 lines across 130 files. The function field portion comprises 17
files and 7{,}372 lines with 364 declarations, all verified against
Mathlib with zero \texttt{sorry} axioms.

\subsection{File inventory}

The 17 files divide into two categories: 10 \textbf{green} files (fully
unconditional, zero open hypotheses) totaling 4{,}073 lines, and 7
\textbf{yellow} files (containing or depending on open propositions)
totaling 3{,}299 lines.

\begin{center}
\renewcommand{\arraystretch}{1.2}
\begin{longtable}{llrl}
\toprule
\textbf{File} & \textbf{Status} & \textbf{Lines} & \textbf{Key contribution} \\
\midrule
\endhead
\texttt{Analog.lean} & Yellow & 952 & FFEMData, FFMullinConjecture, coprimality \\
\texttt{AutonomousMap.lean} & Green & 225 & $\Phi_3$ criterion, unreachability \\
\texttt{Bootstrap.lean} & Yellow & 438 & FFDH $\Rightarrow$ FFMC, injectivity \\
\texttt{CyclicWalkCoverage.lean} & Green & 747 & Weyl criterion, $\ZZ/4\ZZ$ counterex. \\
\texttt{FactorTree.lean} & Yellow & 455 & Mixed selection, walk products \\
\texttt{FFCharacterSums.lean} & Green & 235 & Char.\ cancellation, PNT-in-APs, FCD \\
\texttt{FFSieve.lean} & Green & 327 & \textbf{Headline}: almost-all unconditional \\
\texttt{Finiteness.lean} & Green & 188 & Finite irreds per degree \\
\texttt{IrreducibilityDensity.lean} & Green & 166 & $Q \mid t^{p^n} - t$, density bound \\
\texttt{Master.lean} & Green & 306 & Theorems A/B/C/D assembly \\
\texttt{MultiplierCCSB.lean} & Yellow & 763 & Population CCSB from Weil \\
\texttt{NecklaceFormula.lean} & Green & 446 & \textbf{Necklace identity} ($\sim$380 lines) \\
\texttt{OrbitBarrier.lean} & Green & 300 & Dead Ends \#127/\#129/\#130 \\
\texttt{PopulationEquidist.lean} & Yellow & 370 & Factor pool degree growth \\
\texttt{StochasticMC.lean} & Green & 431 & Stochastic MC, phase transition \\
\texttt{SubgroupEscape.lean} & Yellow & 552 & Weil $\Rightarrow$ SE, threshold \\
\texttt{WeakMC.lean} & Yellow & 471 & Degree escape, pool partition \\
\midrule
\textbf{Total} & & \textbf{7{,}372} & \textbf{364 declarations} \\
\bottomrule
\end{longtable}
\end{center}

\subsection{The necklace identity proof}

The necklace identity (Theorem~\ref{thm:necklace}) is the foundational
result of the sieve chain. Its Lean proof runs approximately 380 lines
and proceeds through six stages:

\begin{enumerate}
  \item \textbf{Divisibility transitivity} (lines 98--106): if $d \mid n$
    and $Q$ is monic irreducible of degree $d$, then $Q \mid t^{p^n} - t$.
    Uses \texttt{monic\_irreducible\_dvd\_X\_pow\_sub\_X} and
    \texttt{dvd\_pow\_pow\_sub\_self\_of\_dvd}.

  \item \textbf{Product divisibility} (lines 119--132): the product of all
    monic irreducibles with degree dividing $n$ divides $t^{p^n} - t$.
    Uses \texttt{Finset.prod\_dvd\_of\_coprime} with the coprimality of
    distinct monic irreducibles.

  \item \textbf{Degree constraint via Galois theory} (lines 140--196): if
    a monic irreducible $Q$ divides $t^{p^n} - t$, then $\deg(Q) \mid n$.
    The proof constructs a root $\alpha$ of $Q$ in $\GF{p}{n}$ (using
    \texttt{Splits.exists\_eval\_eq\_zero}), identifies $Q$ with the minimal
    polynomial of $\alpha$ (via \texttt{minpoly.irreducible} and
    \texttt{Associated} comparison), and applies the tower law:
    $\deg(Q) = [\Fp(\alpha):\Fp]$ divides $[\GF{p}{n}:\Fp] = n$ by
    \texttt{IntermediateField.finrank\_dvd\_of\_le\_right}.

  \item \textbf{Membership in the product} (lines 201--208): combining
    stages 2 and 3 to show that the product accounts for all irreducible
    factors of $t^{p^n} - t$.

  \item \textbf{Degree computation} (lines 221--239): the degree of the
    product equals $\sum_{d \mid n} d \cdot \irredcount(d)$, using
    \texttt{Polynomial.natDegree\_prod} and the disjointness of
    irreducible sets at different degrees.

  \item \textbf{Reverse inequality via root counting} (lines 251--382): the
    key step. Every element $\alpha \in \GF{p}{n}$ is a root of the product
    (since $\alpha$'s minimal polynomial divides the product). Hence the
    product has at least $p^n$ roots (in $\GF{p}{n}$), so its degree is at
    least $p^n$. Combined with the forward inequality (degree $\leq p^n$
    from divisibility), equality follows.
\end{enumerate}

\subsection{Key Mathlib dependencies}

The formalization makes essential use of the following Mathlib components:

\begin{itemize}
  \item \textbf{Galois theory}: \texttt{Polynomial.Gal},
    \texttt{GaloisField.card}, \texttt{GaloisField.finrank},
    \texttt{IntermediateField.adjoin.finrank}.
  \item \textbf{Finite field theory}: \texttt{FiniteField.isSplittingField\_of\_card\_eq},
    \texttt{FiniteField.X\_pow\_card\_pow\_sub\_X\_ne\_zero},
    \texttt{FiniteField.pow\_card} (Frobenius identity).
  \item \textbf{Polynomial arithmetic}: \texttt{Polynomial.natDegree\_prod},
    \texttt{Polynomial.Monic.natDegree\_mul},
    \texttt{monicEquivDegreeLT}, \texttt{degreeLTEquiv}.
  \item \textbf{Minimal polynomials}: \texttt{minpoly.monic},
    \texttt{minpoly.irreducible}, \texttt{minpoly.dvd},
    \texttt{minpoly.aeval}.
  \item \textbf{Group theory}: \texttt{orderOf\_dvd\_card} (Lagrange),
    \texttt{ZMod.card\_units\_eq\_totient},
    \texttt{Nat.totient\_prime}.
  \item \textbf{Set theory}: \texttt{Set.infinite\_range\_of\_injective},
    \texttt{Set.Finite.not\_infinite} (for competitor exhaustion).
\end{itemize}

\subsection{Open propositions}

The 21 open propositions (across 5 yellow files) represent hypotheses
that are \emph{not} used by any of the headline unconditional theorems.
They encode the orbit-specificity barrier:

\begin{itemize}
  \item \texttt{FFDynamicalHitting}: cofinal $-1$ hits for the FF-EM walk.
  \item \texttt{SelectionBiasNeutral}: greedy selection preserves character
    cancellation.
  \item \texttt{WeilBound}: the Weil character sum bound (a theorem, but
    formalizing the algebraic geometry is beyond current Mathlib).
  \item \texttt{FFPopulationEquidist}: population equidistribution from Weil.
  \item \texttt{FFDSL}: the function field deterministic stability lemma.
\end{itemize}

Critically, \texttt{ff\_almost\_all\_unconditional} has zero open hypotheses
in its transitive closure. Despite being defined in a file that imports
yellow files, the theorem depends only on unconditionally proved declarations.

\subsection{Verification and reproduction}

The complete formalization compiles with
\begin{center}
\texttt{lake build}
\end{center}
from the project root. Build time is approximately 15--20 minutes on a
modern machine. The function field module can be checked independently via
\begin{center}
\texttt{lake build EM.FunctionField.Master EM.FunctionField.OrbitBarrier}
\end{center}
which verifies the full dependency chain.
